\section{Results}
\label{sec:results}

We organise the results by question. \Cref{sec:main} compares \method with all baselines and floors on four datasets,
using the timeline answer format. \Cref{sec:transfer} then repeats the key comparisons under the standard one-line
prompt, on a grounding-tuned model and on NExT-GQA.
\Cref{sec:pruners} looks at how published pruners behave on grounding, and \cref{sec:res-coverage,sec:res-contrast}
analyse the two components. \Cref{sec:budget,sec:frames} study the budget and how to spend it. The remaining sections
compare with upper bounds that use the gold interval (\cref{sec:oracles}), analyse BOLT (\cref{sec:res-retrieval}),
list what did not work (\cref{sec:negative}) and examine the errors that remain (\cref{sec:errors}).

\paragraph{Reading the tables.}
Numbers are mIoU, or mIoU\,/\,R@0.5 where two are given. $\Delta$ is a paired difference, and $^*$ marks differences
whose 95\% paired bootstrap interval excludes 0. Tables built from the early S-EMBER phases (\cref{tab:pcd8,tab:cf,%
tab:scope,tab:replay,tab:errors} and \cref{tab:causal}) use an earlier implementation of PCD, the official parser and
batch size 1. On the S-EMBER rows we checked, the two parsers differ by 0.2--0.7 mIoU. The exception is the oracle
selector with PCD, where they differ by 1.5 mIoU because 17 answers are written in minutes:seconds. Batch size changes
no paired result. All other S-EMBER tables use the minutes:seconds-aware parser. In every table the best
value is in \best{bold} and the second best \second{underlined}, per column unless a caption says per row. Floors
and differences ($\Delta$) are not ranked.

%% ================================================================================================= 6.1 main
\subsection{Main results}
\label{sec:main}

\begin{table*}[t]
  \centering
  \small
  \setlength{\tabcolsep}{5pt}
  \caption{\textbf{Grounding with 10\% of the visual tokens} (ActivityNet: 30\%), mIoU\,/\,R@0.5. Grey rows are floors:
  the best video-free constant interval, tuned on the test questions, and the BOLT\,$+$\,Coverage prediction
  without PCD lengthened by a factor fitted on S-EMBER. Coverage rows
  include BOLT, which is skipped on Charades-STA, where uniform frames are already dense. Empty cells: not yet computed. S-EMBER closed intervals: \cref{tab:closed}.}
  \label{tab:main}
  \begin{adjustbox}{max width=\linewidth}\begin{tabular}{@{}lcccc@{}}
    \toprule
    Method & S-EMBER & QVHighlights & ActivityNet (30\%) & Charades-STA \\
    \midrule
    % \floor{Best constant interval}       & \floor{36.9 / 34.7} & \floor{26.8 / 17.6} & \floor{37.9 / 35.8} & \floor{29.7 / 29.6} \\
    % \floor{Length rescaling of R $+$ C}  & \floor{38.3 / 40.0} & \floor{40.2 / 42.4} & \floor{34.0 / 33.0} & \floor{24.0 / 17.6} \\
    % \midrule
    All tokens (no pruning)          & 28.0 / 26.7 & 29.1 / 28.6 & 22.9 / 18.2 & 16.6 / 14.2 \\
    Random tokens                    & 25.6 / 24.6 & 26.9 / 25.0 & 22.9 / 18.6 & 17.4 / 13.8 \\
    FastVID~\cite{fastvid}           & 17.9 / 15.8 & 19.2 / 15.4 &             & 13.3 / 11.6 \\
    VidCom\textsuperscript{2}~\cite{vidcom2} & 20.4 / 18.7 & 19.9 / 18.2 &             & 12.3 / \phantom{0}8.4 \\
    FlashVID~\cite{flashvid}         & 22.7 / 21.5 & 22.8 / 20.4 & 20.5 / 15.4$^\dagger$ & 14.1 / 10.8 \\
    SemVID~\cite{semvid}             & 23.3 / 20.6 & 24.5 / 21.2 & 22.8 / 18.6 & 14.6 / 11.2 \\
    \midrule
    % Random $+$ PCD              & \second{35.0} / 33.9 & \second{37.8} / \second{39.2} & \second{31.4} / \second{29.4} & \best{26.2} / \best{24.0} \\
    % SemVID $+$ PCD              & 34.8 / \second{35.8} & 34.8 / 34.8 & 30.2 / 27.8 & 21.7 / 19.2 \\
    BOLT + Coverage                    & 32.5 / 34.1 & 34.5 / 35.6 & 25.4 / 21.2 & 16.4 / 13.6 \\
    \oursrow BOLT + Coverage + PCD & \best{39.3} / \best{40.8} & \best{48.8} / \best{54.8} & \best{34.7} / \best{32.4} & \second{23.3} / \second{21.2} \\
    \bottomrule
    \multicolumn{5}{@{}l}{\footnotesize $^\dagger$ at 10\%.}
  \end{tabular}\end{adjustbox}
\end{table*}

\begin{table*}[t]
  \centering
  \small
  \setlength{\tabcolsep}{4pt}
  \caption{\textbf{Paired gains of \method in mIoU} ($^*$: 95\% paired bootstrap interval excludes 0).}
  \label{tab:paired}
  \begin{adjustbox}{max width=\linewidth}\begin{tabular}{@{}lcccc@{}}
    \toprule
    \method \vs & S-EMBER & QVHighlights & ActivityNet (30\%) & Charades-STA \\
    \midrule
    All tokens            & \up{$+11.2^*$}            & \up{$+19.7^*$} & \up{$+11.8^*$} & \\
    Random $+$ PCD   & \up{$+4.2^*$}     & \up{$+11.0^*$}  & \up{$+3.3^*$}   & \dn{$-2.9^*$} \\
    SemVID $+$ PCD   & \up{$+4.5^*$}     & \up{$+14.1^*$} & \up{$+4.5^*$}   & \\
    Best constant         & $+2.4$      & \up{$+22.1^*$} & \dn{$-3.1^*$} & \dn{$-6.4^*$} \\
    Length rescaling      & $+1.0$      & \up{$+8.6^*$}   & $+0.8$    & $-0.7$ \\
    \bottomrule
  \end{tabular}\end{adjustbox}
\end{table*}

\Cref{tab:main,tab:paired} lead to three conclusions.

(i) With the timeline prompt, \method nearly doubles the uncompressed model on the least biased benchmark.
On QVHighlights with 10\% of the tokens, \method reaches 54.8 R@0.5 (R@0.7: 30.2). All tokens give 28.6 (18.4), and
random tokens 25.0 (16.2). \method also clears every floor by a wide margin, with $+22.1$ mIoU over the best constant
and $+8.6$ over length rescaling. This comparison holds the prompt fixed. The one-line prompt of TimeLens-Bench gives
the uncompressed model far higher scores, and there the gain of \method is smaller (\cref{sec:transfer}).

(ii) On S-EMBER, the gain holds where the window end is not a shortcut. \method beats all tokens, random
selection and SemVID under the same decoding. Against the floors its pooled gain is small and not significant,
because two thirds of S-EMBER intervals end at the question time. Closed intervals do not offer that shortcut, and on
them \method gains $+7.6^*$ mIoU over the best constant and $+7.5^*$ over length rescaling
(\cref{tab:closed}).

(iii) On ActivityNet and Charades-STA, the timeline prompt is the bottleneck. With this prompt and our
low-resolution copies, the uncompressed model falls far below the constant floor on both datasets (22.9 \vs 37.9 and
16.6 \vs 29.7). With the one-line prompt of TimeLens-Bench, the same model scores 47.1 and 48.5 mIoU
(\cref{sec:transfer}). On ActivityNet, \method improves on all tokens and on every pruner, yet no method reaches the
floor. On Charades-STA, Coverage is worse than random selection (\cref{sec:frames}). We count both datasets among the
limits of the method.

\begin{table}[t]
  \centering
  \small
  \setlength{\tabcolsep}{3pt}
  \caption{\textbf{Grounded QA on S-EMBER} (10\%; all rows from judging session A, \cref{sec:setup}). $\Delta$Acc.\ \vs
  random tokens.}
  \label{tab:qa}
  \begin{adjustbox}{max width=\linewidth}\begin{tabular}{@{}lcccc@{}}
    \toprule
    Method & Acc. & GQ@0.5 & mIoU & $\Delta$Acc. \\
    \midrule
    All tokens             & 19.8 & \phantom{0}8.8 & 28.0 & $+2.5$ \\
    Random tokens          & 17.3 & \phantom{0}7.4 & 25.6 & -- \\
    FastVID                & 15.2 & \phantom{0}4.8 & 17.9 & $-2.1$ \\
    VidCom\textsuperscript{2} & 16.2 & \phantom{0}5.9 & 20.4 & $-1.1$ \\
    FlashVID               & 17.3 & \phantom{0}6.1 & 22.7 & $+0.0$ \\
    SemVID                 & 19.8 & \phantom{0}7.6 & 23.3 & $+2.5$ \\
    Coverage, uniform frames & \best{22.5} & 11.2 & 31.7 & \up{$+5.3^*$} \\
    Random $+$ PCD    & 20.4 & 12.2 & \second{35.0} & $+3.2$ \\
    SemVID $+$ PCD$^\circ$ & 17.9 & \phantom{0}9.9 & 32.3 & $+0.6$ \\
    \oursrow \method                & 21.7 & \second{13.3} & \best{39.3} & \up{$+4.4^*$} \\
    \oursrow \method $+$ coverage gate & \second{22.1} & \best{13.5} & \best{39.3} & \up{$+4.8^*$} \\
    \bottomrule
    \multicolumn{5}{@{}l}{\footnotesize $^\circ$ earlier implementation of PCD.}
  \end{tabular}\end{adjustbox}
\end{table}

\paragraph{Answers.}
\method also improves the answers (\cref{tab:qa}). Its accuracy is 21.7, against 17.3 for random tokens ($+4.4^*$)
and 19.8 for all tokens. GQ@0.5 rises to 13.3, from 7.4 and 8.8 respectively. These numbers come from judging
session A; in session B the same \method run scores 21.9 and 13.5 (\cref{tab:tl8b,tab:oracle}). Coverage alone gives
the best accuracy (22.5), so PCD trades about one point of accuracy for its grounding gain. Differences below about
two points of accuracy are within the decoding noise of the judge setup. \TODO{SemVID $+$ PCD in the same judging session.}

\begin{table}[t]
  \centering
  \small
  \setlength{\tabcolsep}{3pt}
  \caption{\textbf{S-EMBER by question type}, mIoU\,/\,R@0.5 (10\%).}
  \label{tab:types}
  \begin{adjustbox}{max width=\linewidth}\begin{tabular}{@{}lccc@{}}
    \toprule
    Method & Duration (191) & Counting (184) & Location (100) \\
    \midrule
    All tokens            & 37.6 / 40.3 & 30.1 / 25.5 & \phantom{0}5.9 / 3.0 \\
    Random tokens         & 33.2 / 35.1 & 29.6 / 26.1 & \phantom{0}3.5 / 2.0 \\
    BOLT + Coverage              & \second{43.2} / \second{46.6} & 35.3 / 36.4 & \phantom{0}7.1 / 6.0 \\
    Random $+$ PCD   & 41.0 / 40.8 & 41.9 / 40.2 & \best{11.0} / \best{9.0} \\
    SemVID $+$ PCD   & 38.9 / 36.6 & \second{44.6} / \best{50.0} & \phantom{0}8.9 / \second{8.0} \\
    \oursrow BOCC               & \best{47.6} / \best{50.8} & \best{46.9} / \second{48.4} & \phantom{0}\second{9.3} / \second{8.0} \\
    \bottomrule
  \end{tabular}\end{adjustbox}
\end{table}

\paragraph{Question types.}
\method improves duration and counting questions most (\cref{tab:types}). Location questions (``where did I put
\dots'') remain essentially unsolved, with mIoU $\leq 11$ for every method. Part of the reason is the labels. For 62
of the 100 location questions, the annotators disagree on the interval (mean pairwise IoU $<0.6$), which caps what
any method can score.

%% ================================================================================================= 6.2 transfer
\subsection{Standard prompt, a grounding-tuned model, and short videos}
\label{sec:transfer}

The results above use one model and the reasoning-style timeline prompt on every dataset. We now test the method
under the standard one-line grounding prompt of TimeLens-Bench~\cite{timelens}, on a grounding-tuned model
(TimeLens-8B~\cite{timelens}), and on NExT-GQA~\cite{nextgqa}. Unless stated otherwise, paired differences use the
same questions; unpruned runs use batch size 1--2 and pruned runs batch size 8, which alone moves scores by up to
about two points (\cref{sec:setup}).

\paragraph{Our pipeline reproduces published numbers.}
With the TimeLens evaluation script, Qwen3-VL-8B scores 48.5, 59.6 and 47.1 mIoU on the full Charades-STA, QVHighlights
and ActivityNet sets of TimeLens-Bench (published: 48.3, 59.4, 46.8), and TimeLens-8B 55.4, 65.7 and 53.2 (published:
55.2, 65.5, 53.2). On the same queries, our own pipeline without pruning matches the script: $+0.5$ (Charades, 400
queries), $+1.4$ (ActivityNet, 400) and $0.0$ (QVHighlights, 513) mIoU, none significant (\cref{tab:anchors}).

\begin{table}[t]
  \centering
  \small
  \setlength{\tabcolsep}{3pt}
  \caption{\textbf{Answer format} (QVHighlights, 513 TimeLens-Bench queries, Qwen3-VL-8B, 300 frames, 31\% of the
  tokens), mIoU. One-line: the TimeLens-Bench prompt. Timeline: the prompt of all other tables. }
  \label{tab:prompt}
  \begin{adjustbox}{max width=\linewidth}\begin{tabular}{@{}lcccc@{}}
    \toprule
    & \multicolumn{2}{c}{One-line prompt} & \multicolumn{2}{c}{Timeline prompt} \\
    \cmidrule(lr){2-3} \cmidrule(l){4-5}
    Selection & greedy & PCD & greedy & PCD \\
    \midrule
    All tokens      & \second{60.5} & --   & 30.4 & -- \\
    Random tokens   & 58.2 & --   & \best{33.4} & 43.9 \\
    Coverage        & \best{61.7} & 61.6 & \second{30.7} & 43.9 \\
    \bottomrule
  \end{tabular}\end{adjustbox}
\end{table}

\paragraph{The timeline prompt is costly for single-moment grounding.}
\Cref{tab:prompt} changes only the prompt. On QVHighlights queries, asking for a timeline of observations before the
interval costs about 30 mIoU against the one-line prompt ($-30.0^*$ without pruning, $-31.0^*$ for Coverage). PCD
recovers part of the loss ($+13.2^*$ on Coverage, $+10.4^*$ on random tokens) but stays far below the one-line prompt
($-17.8^*$), and under the one-line prompt it has no effect ($-0.1$, n.s.; also $-0.3$ and $-0.5$ with BOLT). With
the timeline prompt, Coverage is not better than random selection on these queries ($-2.8^*$ greedy, $+0.1$ with PCD).
Two consequences follow. First, the QVHighlights, ActivityNet and Charades-STA numbers of \cref{tab:main} are
relative to a format that handicaps the uncompressed model; the gain over ``all tokens'' there (\cref{sec:main})
does not hold against the one-line prompt. Second, PCD corrects time values in reasoning-style answers, where the
model writes several times before the interval; when the answer is a single interval, there is little for it to
correct.

\begin{table}[t]
  \centering
  \small
  \setlength{\tabcolsep}{3pt}
  \caption{\textbf{Standard grounding prompt} (QVHighlights, 513 TimeLens-Bench queries, one-line prompt),
  mIoU\,/\,R@0.5. Top: 300 frames at protocol size; bottom: 64 frames chosen by BOLT. Kept: visual tokens in
  memory. TimeLens input: the TimeLens-Bench input re-sampled to the same number of tokens. Keep ratios count timestamp
  tokens too: 31\% of all tokens is 24\% of the visual tokens.}
  \label{tab:standard}
  \begin{adjustbox}{max width=\linewidth}\begin{tabular}{@{}lccc@{}}
    \toprule
    Method & Kept & Qwen3-VL-8B & TimeLens-8B \\
    \midrule
    All tokens                         & 13{,}650 & \second{60.5} / \second{64.9} & \second{67.1} / \second{73.1} \\
    TimeLens input, 2\,fps low-res     & 4{,}200  & 58.0 / 62.0 & 64.3 / 68.8 \\
    Random tokens (31\%)               & 3{,}271  & 58.2 / 60.8 & 65.3 / 70.4 \\
    SemVID (31\%)                      & 4{,}232  & 58.8 / 62.4 & -- \\
    \oursrow Coverage (31\%)           & 3{,}271  & \best{61.7} / \best{66.7} & \best{67.3} / \best{73.9} \\
    \midrule
    TimeLens input, 1\,fps low-res     & 1{,}125  & 49.5 / 53.0 & 57.5 / 61.8 \\
    BOLT $+$ random (23\%)         & 897      & \second{55.5} / \second{60.8} & -- \\
    \oursrow BOLT $+$ Coverage (23\%) & 897   & \best{58.3} / \best{63.7} & \best{64.3} / \best{71.2} \\
    \bottomrule
  \end{tabular}\end{adjustbox}
\end{table}

\paragraph{Under the standard prompt, Coverage transfers to both models.}
With the one-line prompt (\cref{tab:standard}), Coverage at 24\% of the visual tokens beats random selection on both
models ($+3.5^*$ and $+2.0^*$ mIoU), beats SemVID ($+2.9^*$), and beats the TimeLens input re-sampled to a similar
budget ($+3.7^*$, $+3.0^*$). It matches the uncompressed model ($+1.2$ for Qwen3-VL-8B, with a lower bound just at
zero; $+0.2$ for TimeLens-8B). At 897 tokens, BOLT with Coverage exceeds a 1\,fps TimeLens input of 1{,}125 tokens
by $+8.8^*$ and $+6.8^*$, and BOLT with random tokens by $+2.8^*$. At this budget
the grounding-tuned model does lose accuracy: TimeLens-8B with BOLT and Coverage is $2.8^*$ mIoU below its
uncompressed memory, whereas at 3{,}271 tokens it is level ($+0.2$).

\begin{table}[t]
  \centering
  \small
  \setlength{\tabcolsep}{2.5pt}
  \caption{\textbf{Full TimeLens-Bench sets} (one-line prompt, 31\% of the tokens), mIoU\,/\,R@0.5. All tokens and low-res:
  the TimeLens evaluation script, the latter with a lower resolution at a matched budget (about 3.2k / 4.2k visual
  tokens on Charades-STA / QVHighlights). Paired $\Delta$mIoU below.}
  \label{tab:fullsweep}
  \begin{adjustbox}{max width=\linewidth}\begin{tabular}{@{}lcccc@{}}
    \toprule
    & \multicolumn{2}{c}{Qwen3-VL-8B} & \multicolumn{2}{c}{TimeLens-8B} \\
    \cmidrule(lr){2-3} \cmidrule(l){4-5}
    Method & Charades & QVHighlights & Charades & QVHighlights \\
    \midrule
    All tokens (no pruning)   & \second{48.5} / \best{54.0} & \second{59.6} / \best{64.5} & \best{55.4} / \best{63.1} & \best{65.7} / \best{71.8} \\
    TimeLens input, low-res   & 46.0 / 49.6 & 57.3 / 61.6 & 52.8 / 59.6 & 63.3 / 68.4 \\
    Random tokens (31\%)      & 48.2 / \second{52.3} & 58.1 / \second{62.6} & 54.5 / \second{61.7} & 64.2 / 69.6 \\
    \oursrow Coverage (31\%)  & \best{48.6} / \best{54.0} & \best{59.7} / \best{64.5} & \second{55.2} / \best{63.1} & \second{65.4} / \second{71.0} \\
    \midrule
    Coverage \vs random       & $+0.5^*$ & $+1.6^*$ & $+0.7^*$ & $+1.2^*$ \\
    Coverage \vs low-res      & $+2.6^*$ & $+2.4^*$ & $+2.4^*$ & $+2.1^*$ \\
    Coverage \vs all tokens   & $+0.1$   & $+0.0$   & $-0.2$   & $-0.3$ \\
    \bottomrule
  \end{tabular}\end{adjustbox}
\end{table}

\paragraph{The same holds on the full sets.}
On the complete Charades-STA (3{,}363 queries) and QVHighlights (1{,}541) sets of TimeLens-Bench
(\cref{tab:fullsweep}), Coverage at 31\% of the tokens is lossless on both models ($-0.3$ to $+0.1$ mIoU, n.s.). It
beats random selection in all four cells ($+0.5^*$ to $+1.6^*$) and the TimeLens input reduced to the same budget by
lowering the resolution in all four ($+2.1^*$ to $+2.6^*$). The gap to random is largest on QVHighlights, where
videos are longer; on Charades-STA, 31\% leaves about 3{,}350 of 10{,}800 tokens and random selection is already
close, as on NExT-GQA below. Keeping full resolution and choosing tokens beats spending the same budget on a lower
resolution: random selection also beats the low-res input on Charades-STA ($+2.2^*$ and $+1.7^*$), though not
significantly on QVHighlights.

\begin{table}[t]
  \centering
  \small
  \setlength{\tabcolsep}{2.5pt}
  \caption{\textbf{A grounding-tuned model on S-EMBER} (10\%, timeline prompt), mIoU\,/\,R@0.5 and Acc.\,/\,GQ@0.5
  (one judging session for all rows).}
  \label{tab:tl8b}
  \begin{adjustbox}{max width=\linewidth}\begin{tabular}{@{}lcccc@{}}
    \toprule
    & \multicolumn{2}{c}{Qwen3-VL-8B} & \multicolumn{2}{c}{TimeLens-8B} \\
    \cmidrule(lr){2-3} \cmidrule(l){4-5}
    Method & mIoU / R@0.5 & Acc. / GQ & mIoU / R@0.5 & Acc. / GQ \\
    \midrule
    All tokens (no pruning)   & 28.0 / 26.7 & 19.6 / \phantom{0}8.8 & 33.1 / 34.1 & 20.6 / 12.2 \\
    Random tokens             & 25.6 / 24.6 & 17.1 / \phantom{0}7.2 & 28.2 / 26.7 & 19.6 / \phantom{0}9.5 \\
    Random $+$ PCD            & \second{35.0} / 33.9 & 20.0 / \second{11.8} & \second{38.4} / \second{39.4} & 20.4 / 12.4 \\
    BOLT $+$ Coverage         & 32.5 / \second{34.1} & \best{22.3} / 11.2 & 37.2 / 37.9 & \second{22.9} / \second{13.1} \\
    \oursrow \method          & \best{39.3} / \best{40.8} & \second{21.9} / \best{13.5} & \best{43.9} / \best{46.9} & \best{27.4} / \best{17.5} \\
    \midrule
    \method \vs all tokens    & $+11.2^*$ & $+2.3$ / $+4.6^*$ & $+10.7^*$ & $+6.7^*$ / $+5.3^*$ \\
    \bottomrule
  \end{tabular}\end{adjustbox}
\end{table}

\paragraph{On grounded QA, the full method transfers to a grounding-tuned model.}
S-EMBER answers explain before they ground, so the timeline format is natural there. TimeLens-8B, trained for
grounding, is above Qwen3-VL-8B in every row of \cref{tab:tl8b}, and the ordering of the methods is the same for both
models. With 10\% of the tokens, \method improves on the uncompressed TimeLens-8B on every metric: $+10.7^*$ mIoU,
$+6.7^*$ accuracy and $+5.3^*$ GQ@0.5, and on random selection by $+15.7^*$ mIoU, $+7.8^*$ accuracy and $+8.0^*$
GQ@0.5. Each component helps on its own: PCD adds $+10.2^*$ mIoU to random tokens and $+6.7^*$ to BOLT $+$ Coverage,
and BOLT $+$ Coverage adds $+9.0^*$ to random tokens. With TimeLens-8B, \method answers and grounds better than
Qwen3-VL-8B with the gold window of \cref{tab:oracle} (GQ@0.5 17.5 \vs 15.8).

\begin{table}[t]
  \centering
  \small
  \setlength{\tabcolsep}{3pt}
  \caption{\textbf{NExT-GQA, multiple choice} (test, 1{,}005 questions, 64 frames, two-query protocol: one query for
  the answer, one for the interval; no PCD). 5\% keeps 406 of 8{,}122 tokens.}
  \label{tab:nextgqa}
  \begin{adjustbox}{max width=\linewidth}\begin{tabular}{@{}llcccc@{}}
    \toprule
    Model & Method & Acc@QA & Acc@GQA & mIoP & mIoU \\
    \midrule
    \multirow{9}{*}{Qwen3-VL-8B} & All tokens     & \second{79.7} & \second{38.9} & \best{48.0} & \best{37.7} \\
                                 & Random 30\%    & 78.6 & 38.6 & 47.5 & 37.3 \\
                                 & SemVID 30\%    & 78.9 & 38.5 & 46.2 & 35.8 \\
                                 & Coverage 30\%  & \best{79.8} & \best{39.4} & \second{47.9} & \second{37.6} \\
                                 & Random 10\%    & 75.3 & 34.9 & 45.2 & 35.3 \\
                                 & Coverage 10\%  & 78.4 & 36.1 & 44.6 & 34.8 \\
                                 & Random 5\%     & 68.3 & 27.3 & 40.5 & 30.5 \\
                                 & SemVID 5\%     & 74.0 & 30.0 & 41.0 & 30.0 \\
                                 & Coverage 5\%   & 75.0 & 32.4 & 41.6 & 31.4 \\
    \midrule
    \multirow{5}{*}{TimeLens-8B} & All tokens     & \best{76.7} & \best{41.5} & \best{52.1} & \best{39.8} \\
                                 & Random 10\%    & 72.3 & 36.9 & 48.9 & 37.7 \\
                                 & Coverage 10\%  & \second{75.4} & \second{38.5} & \second{49.3} & \second{37.8} \\
                                 & Random 5\%     & 66.0 & 28.7 & 43.6 & 34.0 \\
                                 & Coverage 5\%   & 71.5 & 33.3 & 46.0 & 35.6 \\
    \bottomrule
  \end{tabular}\end{adjustbox}
\end{table}

\paragraph{On short videos, the gain of Coverage grows as the budget shrinks.}
NExT-GQA videos last about 40\,s, so 64 frames already sample them densely (\cref{tab:nextgqa}). The advantage of
Coverage over random selection grows as the budget shrinks: on Qwen3-VL-8B it adds $+1.2$ (n.s.), $+3.1^*$ and
$+6.8^*$ Acc@QA at 30\%, 10\% and 5\%, and at 5\% also $+5.2^*$ Acc@GQA. On TimeLens-8B the gains are $+3.1^*$ Acc@QA at
10\%, and at 5\% Coverage beats random selection on every metric ($+5.6^*$ Acc@QA, $+4.7^*$ Acc@GQA, $+2.4^*$ mIoP,
$+1.6^*$ mIoU). At 30\% Coverage is lossless ($+0.5$ Acc@GQA, $-0.2$ mIoU, n.s.) and grounds better than SemVID
($+1.8^*$ mIoU). At 5\%, SemVID also answers well (74.0, $+5.8^*$ over random), and Coverage leads it by $+2.5$
Acc@GQA (interval just including 0). When frames are dense, where tokens are kept matters most for the answer and
less for the interval, as on Charades-STA (\cref{sec:frames}). We use the two-query protocol; asking for answer and
interval in one reply, as in prior zero-shot reports~\cite{lover}, lowers Acc@GQA by $16.5^*$ (Qwen3-VL-8B) and $9.2^*$
(TimeLens-8B), mostly because the time values are then written as minutes:seconds or copied from the prompt's example,
so absolute numbers are comparable only within a protocol.

\begin{table}[t]
  \centering
  \small
  \setlength{\tabcolsep}{3pt}
  \caption{\textbf{NExT-GQA, open-ended} (Qwen3-VL-8B, 1{,}005 questions, options removed). Answers are judged
  correct when they match the ground truth in meaning (LLM judge, Video-ChatGPT prompt~\cite{videochatgpt}); Acc@GQA
  also requires IoP\,$\geq$\,0.5. Two-query: the answer prompt plus the TimeLens grounding query. Joint: one
  timeline answer with the interval.}
  \label{tab:nextgqa-open}
  \begin{adjustbox}{max width=\linewidth}\begin{tabular}{@{}llcccc@{}}
    \toprule
    Protocol & Method & Acc@QA & Acc@GQA & mIoP & mIoU \\
    \midrule
    \multirow{2}{*}{Two-query} & All tokens       & \second{33.5} & \best{17.4} & \best{48.0} & \best{37.7} \\
                               & Coverage 10\%    & 33.0 & \second{16.4} & \second{44.6} & \second{34.8} \\
    \midrule
    \multirow{5}{*}{Joint}     & All tokens       & \best{34.1} & 11.0 & 30.4 & 20.8 \\
                               & Random 10\%      & 31.1 & \phantom{0}9.7 & 30.5 & 19.9 \\
                               & SemVID 10\%      & 28.1 & \phantom{0}7.9 & 28.6 & 17.2 \\
                               & Coverage 10\%    & 33.2 & 10.6 & 30.6 & 20.4 \\
                               & Coverage 10\% $+$ PCD & 32.2 & 10.7 & 33.4 & 24.7 \\
    \bottomrule
  \end{tabular}\end{adjustbox}
\end{table}

\paragraph{Open-ended answers.}
Without the options (\cref{tab:nextgqa-open}), answering is far harder (33.5 \vs 79.7 Acc@QA) while grounding is
unchanged, since the grounding query is the same. Our judge accepts the ground-truth answer for all 1{,}005 questions
and agrees with letter correctness on 99.2\% of the multiple-choice answers. Without pruning, Qwen3-VL-8B reaches 17.4
Acc@GQA zero-shot, against 21.4 for TOGA~\cite{toga}, which is trained for open-ended grounded QA and judged with
GPT-3.5. Coverage at 10\% keeps the answers in both protocols ($-0.5$ and $-0.9$ Acc@QA, n.s.), whereas random selection
and SemVID lose them ($-3.0^*$, $-6.1^*$ in the joint protocol; Coverage leads SemVID by $+5.2^*$ Acc@QA and $+2.8^*$
Acc@GQA). The joint timeline answer grounds far worse than the separate query (20.8 \vs 37.7 mIoU): its listed moments
start where the subject is first visible, so 81\% of its intervals start at 0\,s (ground truth: 19\%). PCD raises its
mIoU ($+3.9^*$) but not Acc@GQA, and on the two-query protocol PCD changes neither ($-0.4$ Acc@GQA, n.s.), as with the
one-line prompt on QVHighlights. On NExT-GQA, the contribution is therefore Coverage; PCD helps only where the answer
format lets the prior choose the times. \TODO{two-query Random and SemVID 10\% (being judged); TimeLens-8B open-ended.}


%% ================================================================================================= 6.2 pruners
\subsection{QA pruners on grounding}
\label{sec:pruners}

% Paired difference to random token selection (mIoU), greedy decoding, 10% of the tokens.
\begin{figure}[t]
\centering
\begin{tikzpicture}
\begin{axis}[
    width=\linewidth, height=5.4cm, ybar, bar width=4.6pt,
    ymin=-9.5, ymax=10, ylabel={$\Delta$mIoU vs.\ random tokens}, ylabel style={font=\sffamily\scriptsize},
    symbolic x coords={fast,vid,flash,sem,cov}, xtick=data,
    xticklabels={FastVID, VidCom$^2$, FlashVID, SemVID, \textbf{Coverage (ours)}},
    xticklabel style={font=\sffamily\scriptsize, rotate=18, anchor=north east}, yticklabel style={font=\sffamily\scriptsize},
    axis x line*=bottom, axis y line*=left,
    legend style={font=\sffamily\scriptsize, at={(0.5,1.03)}, anchor=south, draw=none, legend columns=3,
                  /tikz/every even column/.append style={column sep=6pt}},
    enlarge x limits=0.13, ymajorgrids, grid style={black!8},
    extra y ticks={0}, extra y tick style={grid style={black!60, thick}}]
  \addplot[fill=black!55, draw=black!70] coordinates {(fast,-7.6) (vid,-5.2) (flash,-2.9) (sem,-2.3) (cov,6.9)};
  \addplot[fill=ours!75, draw=ours] coordinates {(fast,-7.7) (vid,-7.0) (flash,-4.0) (sem,-2.4) (cov,7.6)};
  \addplot[fill=orange!60, draw=orange!80!black] coordinates {(fast,-4.1) (vid,-5.2) (flash,-3.3) (sem,-2.8) (cov,-1.0)};
  \legend{S-EMBER, QVHighlights, Charades-STA}
\end{axis}
\end{tikzpicture}
\caption{\textbf{QA pruners fall below random selection on grounding.} Paired mIoU difference to random token
selection at the same 10\% budget, greedy decoding. All four published pruners lose on all three datasets.
Coverage gains on the two long-video datasets and is 1.0 mIoU below random selection on 30\,s Charades videos.}
\label{fig:pruners}
\end{figure}


Every published pruner (\cref{fig:pruners}) falls below random token selection on three of the four datasets. On
S-EMBER the losses range from $-2.3$ to $-7.6$ mIoU, all significant. On QVHighlights they range from $-2.4$ to
$-7.7$, with only SemVID not significant, and on Charades-STA from $-2.8$ to $-5.2$, all significant.
On ActivityNet at 10\%, the input-level pruners are above random, which we discuss in \cref{sec:frames}. These gaps
are not seed noise, because random selection varies little across seeds ($26.0 \pm 0.6$ mIoU over three seeds on
S-EMBER). The ordering follows coverage (\cref{tab:retention}). Methods that empty frames most aggressively lose the
most, and those that spread tokens lose the least. Accuracy, by contrast, barely moves (\cref{tab:qa}). FastVID's
accuracy is within noise of random selection ($-2.1$, n.s.), while its mIoU is 7.6 lower. These pruners were designed
to preserve answers, and they do, but they discard the timing.
PCD helps every pruner ($+7$ to $+12$ mIoU) but does not close the gap to \method.

%% ================================================================================================= 6.3 coverage
\subsection{Analysis of Coverage}
\label{sec:res-coverage}

\begin{table}[t]
  \centering
  \small
  \setlength{\tabcolsep}{3pt}
  \caption{\textbf{Component ladder} (mIoU\,/\,R@0.5) with 10\% of tokens;
  PCD: prior-contrastive decoding.}
  \label{tab:ladder}
  \begin{adjustbox}{max width=\linewidth}\begin{tabular}{@{}lccc@{}}
    \toprule
    Configuration & S-EMBER & QVHighlights & ActivityNet \\
    \midrule
    Uniform, random tokens        & 25.6 / 24.6 & 26.9 / 25.0 & 22.9 / 18.6 \\
    BOLT, random tokens              & 29.4 / 29.5 & 28.7 / 28.0 & -- \\
    Coverage (uniform frames)            & 31.7 / 32.4 & 33.7 / 33.0 & 24.6 / \second{21.6} \\
    BOLT $+$ Coverage                       & \second{32.5} / \second{34.1} & \second{34.5} / \second{35.6} & \second{25.4} / 21.2 \\
    \oursrow BOLT $+$ Coverage $+$ PCD & \best{39.3} / \best{40.8} & \best{48.8} / \best{54.8} & \best{34.7} / \best{32.4} \\
    \bottomrule
  \end{tabular}\end{adjustbox}
\end{table}

\paragraph{Contribution.}
\Cref{tab:ladder} adds the components one at a time. On the memory side, Coverage gives the largest gain over random
tokens, $+6.1^*$ mIoU on S-EMBER and $+6.8^*$ on QVHighlights. BOLT adds $+0.8$ to $+0.9$ on top, which is not
significant on any dataset. Once tokens are chosen by the question, the choice of encoded frames matters less.
\TODO{Coverage $+$ PCD without BOLT on all datasets.}

\begin{table}[t]
  \centering
  \small
  \setlength{\tabcolsep}{3pt}
  \caption{\textbf{Design choices of Coverage} (S-EMBER, greedy decoding). Frames: U = uniform.
  $\Delta$ \vs the equal quota with question top-$k$ on the same frames.}
  \label{tab:design}
  \begin{adjustbox}{max width=\linewidth}\begin{tabular}{@{}llcc@{}}
    \toprule
    Token selection & Frames & mIoU & $\Delta$ \\
    \midrule
    \textbf{Equal quota, question top-$k$}       & BOLT & 32.5 & -- \\
    Quota $\propto$ SigLIP relevance (floor 50\%) & BOLT & \best{33.0} & $+0.5$ \\
    Quota $\propto$ SigLIP (floor 25\%)          & BOLT & 32.6 & $+0.1$ \\
    Quota $\propto$ SigLIP $+$ attention         & BOLT & \second{32.9} & $+0.4$ \\
    Quota $\propto$ attention mass               & BOLT & 32.1 & $-0.4$ \\
    % Contrastive token score$^\ddagger$           & BOLT & \best{33.0} & $+0.5$ \\
    One token set shared by all layers           & BOLT & 31.9 & $-0.6$ \\
    \midrule
    \textbf{Equal quota, question top-$k$}       & U & 31.7 & -- \\
    Quota $\propto$ SigLIP relevance             & U & 31.4 & $-0.2$ \\
    Equal quota, sampling instead of top-$k$     & U & 30.0 & $-1.7$ \\
    % Contrastive token score$^\ddagger$           & U & 30.5 & $-1.2$ \\
    % HERMES score, equal quota (stratified)       & U & 23.7 & \\
    % HERMES score, global top-$k$                 & U & 18.3 & \\
    % Question, global top-$k$, no quota (timeline prompt) & U & \TODO{} & \\  -- queued; official-prompt result in tab:official
    \bottomrule
    % \multicolumn{4}{@{}l}{\footnotesize $^\ddagger$ $\log a(\text{question}) - \log a(\text{generic sentence})$.}
  \end{tabular}\end{adjustbox}
\end{table}

\paragraph{Design choices.}
\Cref{tab:design} varies one choice at a time. If there is a quota, an equal one works as well as any weighted one.
We tried five relevance-weighted variants, using SigLIP similarity, attention mass or both, with floors of
25--50\%. All are within $\pm 0.5$ mIoU of the equal quota and none is significant (the uniform-frame SigLIP variant
was our pre-registered primary test). Quotas that follow raw attention mass get pulled by attention sinks and recency
bias, so that a single frame received up to 227 of about 2{,}140 kept tokens. The score itself should come from the
actual question. HERMES's generic probe questions give 23.7 even with an equal quota, and 18.3 without one. Finally,
deterministic top-$k$ beats sampling from a coverage-floored density (31.7 \vs 30.0). A contrastive token score
that subtracts the attention of a generic sentence, meant to cancel sinks, does not help.

\paragraph{Is the quota needed?}
For the generic HERMES score, the equal quota is worth $+5.3^*$ mIoU, because that score concentrates on the end of
the window and empties earlier frames (\cref{tab:retention}). With the actual question's attention, the quota was not
needed in either of the two settings we tested. There, a global top-$k$ of the same per-layer attention, with no
quota, matches the equal quota. On S-EMBER under the official prompt the difference is $-0.1$ mIoU (interval
$[-1.5, +1.3]$, \cref{tab:official}). On QVHighlights with BOLT frames and the one-line prompt it is $-0.0$
($[-0.9, +0.7]$) on the 584 queries run so far, where the global top-$k$ is $+1.8^*$ above random tokens. In these
settings, the gain of Coverage over random selection comes from scoring tokens by the question. The equal quota acts
as a guarantee that every anchor survives, which matters for scores that collapse onto a few frames.
\TODO{global top-$k$ without quota under the timeline prompt on
S-EMBER and on the complete QVHighlights set (queued); share of frames left empty by it.}

\paragraph{Per-layer selection.}
\Cref{tab:shared} compares one token set per layer with a single set shared by all layers. Sharing does not change
random selection, so per-layer diversity is not why random selection beats input-level pruners. For \method, sharing
costs $-1.5^*$ mIoU. Different layers attend to different tokens, and question-driven selection benefits from letting
each layer choose its own.

\begin{table}[t]
  \centering
  \small
  \setlength{\tabcolsep}{3.5pt}
  \caption{\textbf{Per-layer \vs shared token sets} (S-EMBER, 10\%). $\Delta$ = shared $-$ per layer.}
  \label{tab:shared}
  \begin{adjustbox}{max width=\linewidth}\begin{tabular}{@{}lccc@{}}
    \toprule
    Method & Per layer & Shared & $\Delta$mIoU \\
    \midrule
    Random tokens           & 25.6 & 26.3 & $+0.7$ \\
    Random $+$ PCD     & \second{35.0} & \second{35.8} & $+0.7$ \\
    BOLT $+$ Coverage                 & 32.5 & 31.9 & $-0.6$ \\
    \oursrow BOCC                 & \best{39.3} & \best{37.8} & $-1.5$ \\
    \bottomrule
  \end{tabular}\end{adjustbox}
\end{table}

%% ================================================================================================= 6.4 contrast
\subsection{Analysis of PCD}
\label{sec:res-contrast}

\begin{table*}[t]
  \centering
  \small
  \setlength{\tabcolsep}{5pt}
  \caption{\textbf{PCD on eight memories} (S-EMBER, Acc.\,/\,mIoU\,/\,GQ@0.5; earlier implementation of PCD). Plain PCD: on all
  tokens with constant $\alpha{=}0.5$; PCD: confidence-adaptive strength on time values only (ours). $\Delta$ PCD
  \vs greedy. Best and second best per row.}
  \label{tab:pcd8}
  \begin{adjustbox}{max width=\linewidth}\begin{tabular}{@{}lcccccc@{}}
    \toprule
    Memory & Greedy & Plain PCD & PCD & $\Delta$mIoU & $\Delta$GQ@0.5 & Early starts \\
    \midrule
    Random 5\%            & \second{15.8} / 23.6 / 5.9  & 15.4 / \second{29.2} / \second{9.1}  & \best{16.4} / \best{31.5} / \best{9.9}  & \up{$+8.0^*$} & \up{$+4.0^*$} & 39\% $\rightarrow$ 20\% \\
    Random 10\%           & 15.2 / 25.8 / 7.2  & \best{16.2} / \second{32.1} / \best{10.1} & \best{16.2} / \best{32.6} / \best{10.1} & \up{$+6.8^*$} & $+2.9$   & 38\% $\rightarrow$ 19\% \\
    Random 25\%           & 17.7 / 28.2 / 8.2  & \best{18.9} / \second{33.0} / \second{11.8} & \best{18.9} / \best{34.9} / \best{13.3} & \up{$+6.7^*$} & \up{$+5.1^*$} & 37\% $\rightarrow$ 21\% \\
    All tokens            & 18.5 / 28.0 / 8.6  & \second{19.8} / \second{31.5} / \second{11.2} & \best{21.5} / \best{32.7} / \best{13.3} & \up{$+4.6^*$} & \up{$+4.6^*$} & 41\% $\rightarrow$ 23\% \\
    HERMES 10\%           & 13.1 / 18.3 / 4.0  & \best{15.4} / \second{22.6} / \best{7.2}  & \second{13.3} / \best{24.1} / \second{5.9}  & \up{$+5.7^*$} & $+1.9$   & 56\% $\rightarrow$ 31\% \\
    Stratified 10\%       & 14.1 / 23.7 / 6.1  & \second{14.5} / \second{29.9} / \second{8.8}  & \best{15.8} / \best{31.0} / \best{10.5} & \up{$+7.3^*$} & \up{$+4.4^*$} & 36\% $\rightarrow$ 17\% \\
    HERMES streaming 4k   & 13.5 / 21.6 / 4.0  & \best{17.1} / \second{27.4} / \best{8.2}  & \second{15.4} / \best{28.6} / \best{8.2}  & \up{$+6.9^*$} & \up{$+4.2^*$} & 45\% $\rightarrow$ 24\% \\
    HERMES streaming 6k   & 14.3 / 22.6 / 4.8  & \second{14.7} / \second{28.1} / \second{7.6}  & \best{15.4} / \best{29.8} / \best{8.4}  & \up{$+7.2^*$} & \up{$+3.6^*$} & 46\% $\rightarrow$ 24\% \\
    \midrule
    Mean of 8             & 15.3 / 24.0 / 6.1  & \second{16.5} / \second{29.2} / \second{9.3}  & \best{16.6} / \best{30.7} / \best{10.0} & $+6.7$   & $+3.9$   & \\
    \bottomrule
  \end{tabular}\end{adjustbox}
\end{table*}

\paragraph{PCD on other memories.}
With reasoning-style answers, PCD can be added to any memory. \Cref{tab:pcd8} applies it to eight: three random
budgets, the uncompressed model, HERMES and stratified pruning, and two HERMES streaming memories. PCD improves mIoU
significantly on all eight ($+4.6$ to $+8.0$) and GQ@0.5 on six. It roughly halves the early starts everywhere, and
its gains are largest where the memory is smallest. With 10\% of the tokens and random selection, it already beats
the uncompressed model (32.6 \vs 28.0). PCD also transfers to FastVID, the pruner that loses the most coverage, where
it adds $+7.6^*$ mIoU and cuts the leak from $38\% \rightarrow 12\%$.

\begin{table}[t]
  \centering
  \small
  \setlength{\tabcolsep}{3pt}
  \caption{\textbf{What to contrast against} (S-EMBER, random 10\%). 
  % Only the counterfactual input changes; the main input is always the real memory.
  }
  \label{tab:cf}
  \begin{adjustbox}{max width=\linewidth}\begin{tabular}{@{}lccc@{}}
    \toprule
    Counterfactual & mIoU & GQ@0.5 & Early \\
    \midrule
    None (greedy)                         & 25.8 & \phantom{0}7.2 & 38\% \\
    \textbf{No video (ours)}              & \best{32.1} (\up{$+6.2^*$}) & \best{10.1} (\up{$+2.9^*$}) & \best{19\%} \\
    Noised frames (VCD~\cite{vcd})        & 26.3 ($+0.5$) & \phantom{0}\second{7.8} ($+0.6$) & \second{30\%} \\
    % Permuted timestamps                   & $+0.9$ (n.s.) & -- & -- \\
    \midrule
    Evidence-against rule, $\alpha{=}1$\,/\,2 & \second{26.5}\,/\,27.3 & -- & -- \\
    Attention gain (PAI-style)            & $+1.3$ (n.s.) & $+0.8$ & 42\% \\
    Instruction ``do not assume early''   & 25.7 & \phantom{0}6.9 & 37\% \\
    \bottomrule
  \end{tabular}\end{adjustbox}
\end{table}

\paragraph{Choice of counterfactual.}
\Cref{tab:cf} changes only what the memory is contrasted against. The no-video input works best because it isolates
the prior, and contrasting against it removes the prior (early starts $38\% \rightarrow 19\%$). Noised frames, as in
VCD, still show part of the video \emph{and} still produce the prior. The two runs then disagree mainly about content,
and little of the prior is removed ($+0.5$ mIoU). Permuted timestamps also leave the prior in place. A one-sided rule
that only penalises tokens the video argues against is worse than PCD ($-5.6^*$ and $-4.8^*$). Without video the
model is about 99\% sure of ``0'', so even a well-supported ``0'' looks like evidence against. Boosting visual
attention during decoding, or stating the bias in the instruction, does not remove the prior either.

\begin{table}[t]
  \centering
  \small
  \setlength{\tabcolsep}{3pt}
  \caption{\textbf{Time scope and adaptive strength} (S-EMBER, Acc.\,/\,mIoU\,/\,GQ@0.5). Top: PCD on all tokens
  \vs time values only. Bottom: the 78 questions whose evidence truly starts in the first 5\% of the window.
  Top: best per row.}
  \label{tab:scope}
  \begin{adjustbox}{max width=\linewidth}\begin{tabular}{@{}lcc@{}}
    \toprule
    Memory & All tokens & Time values only \\
    \midrule
    Random 10\%  & \best{16.2} / 32.1 / 10.1 & 16.0 / 32.1 / \best{10.3} \\
    Random 25\%  & \best{18.9} / 33.0 / 11.8 & 17.7 / \best{34.0} / 11.8 \\
    All tokens   & 19.8 / 31.5 / 11.2 & \best{20.8} / \best{32.9} / \best{12.4} \\
    HERMES 10\%  & \best{15.4} / 22.6 / \phantom{0}7.2 & 14.7 / \best{24.3} / \phantom{0}\best{7.4} \\
    Streaming 4k & \best{17.1} / 27.4 / \phantom{0}8.2 & 16.2 / \best{29.2} / \phantom{0}\best{9.3} \\
    \midrule
    \multicolumn{3}{@{}l}{\emph{Early-evidence questions (78), all tokens memory}} \\
    Greedy                & \multicolumn{2}{c}{\second{21.8} / 49.2 / \second{16.7}} \\
    PCD, constant $\alpha$ & \multicolumn{2}{c}{16.7 / \second{55.3} / 12.8} \\
    PCD, adaptive $\alpha$ & \multicolumn{2}{c}{\best{23.1} / \best{56.9} / \best{19.2}} \\
    \bottomrule
  \end{tabular}\end{adjustbox}
\end{table}

\paragraph{Time scope and adaptive strength.}
Contrasting every token also changes words that have nothing to do with time. The answer text changed in 85\% of
questions, and lists of more than eight moments rose from 12 to 56. Over-counting rose from 14\% to 34\% of counting
questions. Restricting PCD to time values keeps the grounding gain and disturbs the text less
(\cref{tab:scope}, top). PCD also risks hurting questions whose evidence really starts at 0\,s. With a constant
strength, their accuracy on the uncompressed memory falls from 21.8 to 16.7. The adaptive strength fades when the
model is confident, and with it accuracy rises to 23.1 instead (\cref{tab:scope}, bottom; $+6.4^*$ \vs constant).
The adaptive rule also gives the best \emph{discrimination}, defined as early starts on early-evidence questions minus
early starts on the others. On the all-tokens memory it reaches 56.4, against 49.6 for greedy decoding.

\begin{table}[t]
  \centering
  \small
  \setlength{\tabcolsep}{3pt}
  \caption{\textbf{Decision rules replayed at the first time value} (S-EMBER, random 10\%). Share of answers starting
  at ``0''.}
  \label{tab:replay}
  \begin{adjustbox}{max width=\linewidth}\begin{tabular}{@{}lcc@{}}
    \toprule
    Rule & Early evidence $\uparrow$ & Other (leak) $\downarrow$ \\
    \midrule
    Greedy                       & \best{69\%} & 16\% \\
    PCD, $\alpha{=}0.5$          & 46\% & \phantom{0}\best{6\%} \\
    PCD, adaptive                & \second{51\%} & \phantom{0}\second{7\%} \\
    Evidence-against, $\alpha{=}4$ & \second{51\%} & \phantom{0}9\% \\
    Skewed or tempered prior     & 50--55\% & 8--11\% \\
    Best threshold (uses gold)   & 46\% & \phantom{0}\best{6\%} \\
    \bottomrule
  \end{tabular}\end{adjustbox}
\end{table}

\paragraph{Other decision rules.}
We logged both distributions at the first time value of every answer and replayed other decision rules offline
(\cref{tab:replay}). Even a threshold tuned on the gold labels cannot beat PCD's
trade-off between keeping true early starts and removing leaked ones, so PCD is close to what decoding alone can do.
In the early-start errors that remain, the memory itself is unsure, and fixing them would take more evidence in the
memory.

\paragraph{Is PCD more than a longer interval?}
\label{sec:lengthening}
PCD lengthens predicted intervals (median predicted/gold length $0.57 \rightarrow 0.67$ on S-EMBER), and longer
intervals alone raise mIoU on benchmarks with long gold intervals. The length-rescaling floor tests this directly
(\cref{tab:paired}, last row). \method clearly beats rescaling on QVHighlights ($+8.6^*$) and on S-EMBER closed
intervals ($+7.5^*$, \cref{tab:closed}). On these sets PCD moves the interval toward the evidence instead of only
stretching it. On pooled S-EMBER, ActivityNet and Charades-STA the two are statistically tied, so there the benefit
of PCD cannot be separated from a calibrated lengthening. \TODO{start/end error decomposition on closed questions; mismatched-video and shuffled-frame
negatives.}

\begin{table}[t]
  \centering
  \small
  \setlength{\tabcolsep}{3pt}
  \caption{\textbf{Official S-EMBER prompt} (direct answer with \texttt{Answer:} and \texttt{Time:} lines; Qwen3-VL-8B,
  64 frames, 10\%, second cluster; grounding only, answers not yet judged). mIoU\,/\,R@0.5 with the
  minutes:seconds-aware parser; mIoU with the official parser and the number of answers it cannot parse.}
  \label{tab:official}
  \begin{adjustbox}{max width=\linewidth}\begin{tabular}{@{}lccc@{}}
    \toprule
    Method & mIoU / R@0.5 & Official mIoU & Unparsed \\
    \midrule
    All tokens                          & \TODO{queued} & & \\
    Random tokens                       & 27.9 / 26.1 & 26.4 & \second{17} \\
    Coverage, uniform frames            & \second{30.1} / \second{29.3} & \second{28.8} & \best{13} \\
    Global question top-$k$, no quota   & 30.0 / \second{29.3} & 28.5 & \second{17} \\
    \oursrow \method                    & \best{34.6} / \best{34.3} & \best{29.5} & 73 \\
    \bottomrule
  \end{tabular}\end{adjustbox}
\end{table}

\paragraph{Prompt format.}
We do not yet have a clean PCD-only comparison on S-EMBER under the official direct-answer prompt. An earlier run
that was not matched on the other settings gave $+0.4$ mIoU. Under this prompt (\cref{tab:official}), \method is
$+4.5^*$ mIoU above Coverage alone and $+6.8^*$ above random tokens. This gain mixes PCD with BOLT, since on the
first 116 questions BOLT $+$ Coverage without PCD already gains $+5.3^*$ over Coverage. \method is also $4.7^*$ below
\method with the timeline prompt. Coverage is $+2.2$ (n.s.) above random tokens. PCD makes the model write more times
as minutes:seconds, and 73 \method answers cannot be parsed by the official parser, against 13--17 for the greedy
rows. For the same reason, PCD alone lowers mIoU on NExT-GQA under the same prompt ($-1.5^*$ and $-2.5^*$,
\cref{sec:transfer}). Until the uncompressed model has been run under this prompt, we cannot say whether \method
improves on it. Compared with the uncompressed model in the timeline format, \method with the official prompt is
$+6.6^*$ and Coverage $+2.1$ (n.s.). \TODO{uncompressed model under the official prompt (queued); judged answers; PCD with
and without BOLT under this prompt.}
We also tried an \emph{event} timeline that lists start and end per occurrence (\texttt{Event: 12.0 -- 15.5 seconds}).
Because it states extents, it improves the greedy baseline by $+7.4^*$ mIoU. PCD still adds $+2.1$ to $+3.3$ on top,
reaching 35.0 mIoU and 12.6 GQ@0.5 at 768 tokens. The two fixes overlap. The format corrects extents, and PCD
corrects the leak. \method uses the simpler timeline in all main results. \TODO{event timeline with \method.}

\paragraph{When does PCD help?}
The results so far suggest a simple rule. PCD helps in proportion to how much the video-free prior decides the time
that is written. Its gain is largest where the memory says least, $+8.0^*$ mIoU for random 5\% against $+4.6^*$ for
the uncompressed model (\cref{tab:pcd8}). The adaptive strength fades when the memory makes the model confident, and
so protects true early starts that a constant strength removes (\cref{tab:scope}). With a one-line answer, where the
model states a single interval instead of listing observations, PCD has no effect (\cref{sec:transfer}; those runs
had no answer reserve). One qualification is needed, because PCD can also change the \emph{form} of the answer. On
NExT-GQA it lowers TimeLens-8B's grounding in the timeline format ($-4.4^*$ mIoU) by producing more replies without a
parsable interval. Under the official prompt it makes Qwen3-VL-8B write more times as minutes:seconds
(\cref{sec:transfer}). In short, PCD helps when the prior chooses the times and the answer keeps its format. We state
this as a hypothesis, which the per-token log-probabilities we already store can test directly. \TODO{PCD gain against the divergence between
$p_\Mem$ and $p_\varnothing$ at each written time value, across all conditions; one prediction stated before the run.}

%% ================================================================================================= 6.5 budget
\subsection{Budget}
\label{sec:budget}

\begin{table*}[t]
  \centering
  \small
  \setlength{\tabcolsep}{4.2pt}
  \caption{\textbf{Budget sweep at 64 frames} (mIoU). The second row gives the number of kept tokens per memory.}
  \label{tab:budget}
  \begin{adjustbox}{max width=\linewidth}\begin{tabular}{@{}lccccccccc@{}}
    \toprule
    & \multicolumn{3}{c}{S-EMBER} & \multicolumn{3}{c}{QVHighlights} & \multicolumn{3}{c}{ActivityNet} \\
    \cmidrule(lr){2-4} \cmidrule(lr){5-7} \cmidrule(l){8-10}
    Budget & 5\% & 10\% & 20\% & 5\% & 10\% & 20\% & 10\% & 30\% & 50\% \\
    Kept tokens & ${\sim}$1{,}070 & ${\sim}$2{,}140 & ${\sim}$4{,}280 & ${\sim}$260 & ${\sim}$520 & ${\sim}$1{,}040 & 120--300 & 400--900 & 600--1{,}500 \\
    \midrule
    Random                  & 23.9 & 25.6 & 27.4 &      & 26.9 &      & 17.8 & 22.9 & 23.5 \\
    SemVID                  & 20.9 & 23.3 & 26.5 &      & 24.5 &      & 21.2 & 22.8 & 22.3 \\
    R $+$ C (no PCD)   & \second{34.6} & 32.5 & 31.0 &      & 34.5 &      & 11.3 & 25.4 & 25.0 \\
    SemVID $+$ PCD     & 31.9 & 34.8 & 35.3 &      & 34.8 &      & \best{26.8} & 30.2 & 31.4 \\
    Random $+$ PCD     & 34.5 & \second{35.0} & \second{36.5} &      & \second{37.8} &      & \second{21.5} & \second{31.4} & \second{32.2} \\
    \oursrow \method                 & \best{39.3} & \best{39.3} & \best{36.8} & 9.7 & \best{48.8} & 49.4 & 11.6 & \best{34.7} & \best{34.4} \\
    \bottomrule
  \end{tabular}\end{adjustbox}
\end{table*}

% (a) S-EMBER mIoU vs budget at 64 frames. (b, c) Fixed token budget spent on 64 / 32 / 16 frames.
\begin{figure*}[t]
\centering
\resizebox{\linewidth}{!}{\begin{tikzpicture}
\begin{groupplot}[
    group style={group size=3 by 1, horizontal sep=1.6cm},
    width=0.35\linewidth, height=5.2cm,
    xlabel style={font=\sffamily\scriptsize}, ylabel style={font=\sffamily\scriptsize},
    ticklabel style={font=\sffamily\scriptsize}, title style={font=\sffamily\small},
    ymajorgrids, grid style={black!8}, axis x line*=bottom, axis y line*=left,
    every axis plot/.append style={thick, mark size=2.2pt}]
%% (a) budget
\nextgroupplot[title={(a) S-EMBER: budget at 64 frames}, xlabel={kept tokens (\% of 64 frames)}, ylabel={mIoU},
    xtick={5,10,20}, xmin=3.5, xmax=21.5, ymin=18, ymax=42,
    legend style={font=\sffamily\scriptsize, at={(0.0,-0.3)}, anchor=north west, legend columns=6, draw=none,
                  /tikz/every even column/.append style={column sep=10pt}},
    legend cell align=left]
  \addplot[black!45, mark=o] coordinates {(5,23.9) (10,25.6) (20,27.4)};             \addlegendentry{random}
  \addplot[black!45, dashed, mark=square] coordinates {(5,20.9) (10,23.3) (20,26.5)}; \addlegendentry{SemVID}
  \addplot[orange!80!black, mark=triangle] coordinates {(5,34.6) (10,32.5) (20,31.0)}; \addlegendentry{R + Coverage}
  \addplot[black!70, mark=*] coordinates {(5,34.5) (10,35.0) (20,36.5)};              \addlegendentry{random + Contrast}
  \addplot[black!70, dashed, mark=square*] coordinates {(5,31.9) (10,34.8) (20,35.3)};\addlegendentry{SemVID + Contrast}
  \addplot[ours, very thick, mark=diamond*, mark size=3pt] coordinates {(5,39.3) (10,39.3) (20,36.8)}; \addlegendentry{\textbf{\method}}
  \addplot[floorc, densely dotted, thick, mark=none, forget plot] coordinates {(3.5,36.9) (21.5,36.9)};
  \node[font=\sffamily\tiny, text=floorc, anchor=south west] at (axis cs:3.7,36.9) {best constant};
%% (b) QVH frames
\nextgroupplot[title={(b) QVHighlights: $\sim$260 tokens}, xlabel={frames (same kept tokens)}, ylabel={mIoU},
    xtick={16,32,64}, xmin=10, xmax=70, ymin=0, ymax=50]
  \addplot[black!70, mark=*] coordinates {(16,33.3) (32,32.4)};
  \addplot[black!70, dashed, mark=square*] coordinates {(16,30.6) (32,35.5)};
  \addplot[ours, very thick, mark=diamond*, mark size=3pt] coordinates {(16,40.4) (32,43.1) (64,9.7)};
  \addplot[floorc, densely dotted, thick, mark=none] coordinates {(10,26.8) (70,26.8)};
  \node[font=\sffamily\tiny, text=ours, anchor=east, align=right] at (axis cs:60,12)
    {collapse: 2--4 tokens\\per frame pair};
%% (c) ANet frames
\nextgroupplot[title={(c) ActivityNet: 120--300 tokens}, xlabel={frames (same kept tokens)}, ylabel={mIoU},
    xtick={16,32,64}, xmin=10, xmax=70, ymin=0, ymax=50]
  \addplot[black!70, mark=*] coordinates {(16,30.5) (32,29.5) (64,21.5)};
  \addplot[black!70, dashed, mark=square*] coordinates {(16,30.0) (32,28.4) (64,26.8)};
  \addplot[ours, very thick, mark=diamond*, mark size=3pt] coordinates {(16,32.0) (32,33.9) (64,11.6)};
  \addplot[floorc, densely dotted, thick, mark=none] coordinates {(10,37.9) (70,37.9)};
\end{groupplot}
\end{tikzpicture}}
\caption{\textbf{Budget and how to spend it.} \textbf{(a)}~At 64 frames, question-driven selection improves as the
budget shrinks while random selection and SemVID degrade; \method leads at every budget. \textbf{(b, c)}~At a fixed,
very small number of kept tokens, 64 frames leave 2--4 tokens per frame pair and \method collapses; spending the same
tokens on 32 frames recovers it (QVHighlights 9.7\,$\rightarrow$\,43.1, ActivityNet 11.6\,$\rightarrow$\,33.9). Dotted
grey: best video-free constant.}
\label{fig:budget}
\end{figure*}


\Cref{tab:budget,fig:budget} vary the budget at 64 frames and show two regimes. Above roughly 500 kept tokens,
question-driven selection stays flat or even improves as the budget \emph{shrinks}. On S-EMBER, BOLT $+$ Coverage
rises from 31.0 at 20\% to 34.6 at 5\%, and \method from 36.8 to 39.3. Random selection and SemVID degrade over the
same range. Pruning by the question thus removes distractors as well as cost. At 5\%, \method leads SemVID $+$ PCD
by $+7.4^*$ mIoU, and at 20\% by $+1.5$ (n.s.).

Below roughly 300 kept tokens (QVHighlights 5\%, ActivityNet 10\%), the equal quota leaves only 2--4 tokens per frame
pair. The top-attended tokens then carry little content, and \method collapses to 9.7 and 11.6. This is below random
selection and below SemVID, whose cross-frame merging degrades more gracefully. What matters is the absolute number
of tokens per frame, which is low in our re-encoded QVHighlights and ActivityNet copies. One caveat applies. These runs
count the kept timestamps inside the budget under the earlier timestamp rule, which on NExT-GQA left some low-budget
questions (about 280 tokens) with no visual tokens (\cref{sec:transfer}). Part of the collapse below about 300 tokens
may therefore come from this accounting rather than from the equal quota.
\TODO{QVHighlights 5\% and 10\% at 64 frames with the fixed timestamp rule (queued).}

%% ================================================================================================= 6.6 frames vs tokens
\subsection{Fewer frames or fewer tokens per frame?}
\label{sec:frames}

\begin{table*}[t]
  \centering
  \small
  \setlength{\tabcolsep}{5pt}
  \caption{\textbf{Spending a fixed token budget on fewer frames} (mIoU\,/\,R@0.5, all with PCD).
  Within each block every row keeps the same number of tokens: $F$ frames at keep ratio $\rho \cdot 64/F$. Best and second best per row.}
  \label{tab:frames}
  \begin{adjustbox}{max width=\linewidth}\begin{tabular}{@{}llccc@{}}
    \toprule
    Dataset (kept tokens) & Frames $\times$ keep ratio & Random $+$ PCD & SemVID $+$ PCD & \method \\
    \midrule
    \multirow{3}{*}{QVHighlights (${\sim}$260)} & 64 $\times$ 5\%  & --          & --          & \phantom{0}9.7 / \phantom{0}6.6 \\
                                                & 32 $\times$ 10\% & 32.4 / 29.2 & \second{35.5} / \second{33.6} & \best{43.1} / \best{46.6} \\
                                                & 16 $\times$ 20\% & \second{33.3} / \second{32.4} & 30.6 / 28.0 & \best{40.4} / \best{41.0} \\
    \midrule
    \multirow{3}{*}{ActivityNet (120--300)}     & 64 $\times$ 10\% & \second{21.5} / \second{18.2} & \best{26.8} / \best{23.6} & 11.6 / \phantom{0}7.6 \\
                                                & 32 $\times$ 20\% & \second{29.5} / \second{26.8} & 28.4 / 25.8 & \best{33.9} / \best{32.0} \\
                                                & 16 $\times$ 40\% & \second{30.5} / 26.8 & 30.0 / \second{27.6} & \best{32.0} / \best{29.8} \\
    \midrule
    \multirow{2}{*}{Charades-STA (${\sim}$10\% of 64)} & 64 $\times$ 10\% & \best{26.2} / \best{24.0} & 21.7 / 19.2 & \second{23.3} / \second{21.2} \\
                                                       & 32 $\times$ 20\% & \best{24.6} / \best{22.0} & 20.8 / 16.0 & \second{23.6} / \second{21.0}$^\S$ \\
    \bottomrule
    \multicolumn{5}{@{}l}{\footnotesize $^\S$ Coverage on uniform frames; BOLT is skipped on Charades-STA.}
  \end{tabular}\end{adjustbox}
\end{table*}

The collapse at small budgets comes from how the budget is spent, and the equal quota is not at fault in itself. The
timestamp-accounting caveat of \cref{sec:budget} still applies here, since fewer frames also means fewer timestamps
inside the budget. In \cref{tab:frames}, every row of a block keeps the same number of tokens, and only the split
between frames and tokens per frame changes. With 32 frames instead of 64, each frame pair keeps twice as many tokens,
and \method recovers. On QVHighlights it goes from 9.7 to 43.1 mIoU. On ActivityNet it goes from 11.6 to 33.9, which
beats SemVID $+$ PCD by $+5.5$ and nearly matches its own 64-frame result at three times the budget (34.7). Halving
again to 16 frames loses a little, because past a point temporal coverage becomes the bottleneck. The baselines also
gain from fewer frames, but less. On Charades-STA, whose videos last 30\,s, fewer frames do not help any method. We
therefore recommend the rule of \cref{sec:coverage}. Keep at least about eight tokens per frame pair, and reduce $F$
when the budget does not allow it. At equal tokens, more frames with fewer tokens each do not help
either. With 128 frames at 5\% on
S-EMBER, the change is $+1.2$ mIoU (n.s.) without PCD and $-3.6^*$ with it, with 74--78 unparseable answers.

%% ================================================================================================= 6.7 oracles
\subsection{Comparison with gold-informed upper bounds}
\label{sec:oracles}

\begin{table}[t]
  \centering
  \small
  \setlength{\tabcolsep}{2.5pt}
  \caption{\textbf{Upper bounds that use the gold interval} (S-EMBER, Qwen3-VL-8B, 475 questions; Acc.\ and GQ@0.5
  from judging session B; minutes:seconds-aware parser). Oracle selector: the gold frames' tokens are kept first. Gold
  window: all 64 frames inside the gold interval (batch size 3). Leak: intervals starting in the first 5\% of the
  window, on the 397 questions whose evidence starts later.}
  \label{tab:oracle}
  \begin{adjustbox}{max width=\linewidth}\begin{tabular}{@{}lcccccc@{}}
    \toprule
    Memory & Gold & Acc. & mIoU & R@0.5 & GQ@0.5 & Leak \\
    \midrule
    Random 10\%                       & -- & 17.1 & 25.6 & 24.6 & \phantom{0}7.2 & 29\% \\
    Random 10\% $+$ PCD               & -- & 20.0 & 35.0 & 33.9 & 12.0 & \second{14\%} \\
    Oracle selector 10\%              & \checkmark & 18.9 & 27.9 & 25.7 & \phantom{0}9.3 & 33\% \\
    Oracle selector 10\% $+$ PCD      & \checkmark & 18.5 & 34.9 & 35.2 & \phantom{0}9.5 & 16\% \\
    \oursrow \method 10\%             & -- & 21.9 & \second{39.3} & \second{40.8} & 13.5 & \best{11\%} \\
    \midrule
    All tokens                        & -- & 19.6 & 28.0 & 26.7 & \phantom{0}8.8 & 33\% \\
    Gold window                       & \checkmark & \best{29.5} & 36.1 & 38.1 & \best{15.8} & 45\% \\
    Gold window $+$ PCD               & \checkmark & \second{27.2} & \best{41.5} & \best{43.2} & \second{14.9} & 35\% \\
    \bottomrule
  \end{tabular}\end{adjustbox}
\end{table}

\begin{table}[t]
  \centering
  \small
  \setlength{\tabcolsep}{3pt}
  \caption{\textbf{Gold-informed upper bounds for TimeLens-8B} (S-EMBER, 475 questions; all Acc.\ and GQ@0.5 from
  judging session C, \cref{sec:setup}). Settings as in \cref{tab:oracle}; all tokens at batch size 1. No oracle
  selector without PCD was run.}
  \label{tab:oracle-tl8b}
  \begin{adjustbox}{max width=\linewidth}\begin{tabular}{@{}lccccc@{}}
    \toprule
    Memory & Gold & Acc. & mIoU & R@0.5 & GQ@0.5 \\
    \midrule
    Random 10\%                       & -- & 18.5 & 28.2 & 26.7 & \phantom{0}9.1 \\
    Random 10\% $+$ PCD               & -- & 19.8 & 38.4 & 39.4 & 11.8 \\
    Oracle selector 10\% $+$ PCD      & \checkmark & 21.7 & 43.3 & 45.3 & 14.3 \\
    \oursrow \method 10\%             & -- & 26.5 & 43.9 & \second{46.9} & 17.1 \\
    \midrule
    All tokens                        & -- & 19.8 & 33.1 & 34.1 & 11.6 \\
    Gold window                       & \checkmark & \second{28.2} & \second{46.6} & 46.7 & \second{19.2} \\
    Gold window $+$ PCD               & \checkmark & \best{36.2} & \best{54.9} & \best{57.1} & \best{25.9} \\
    \bottomrule
  \end{tabular}\end{adjustbox}
\end{table}

\Cref{tab:oracle} gives the memory information it should not have. On Qwen3-VL-8B, \method beats perfect token
selection at the same budget without using any gold information. Over the oracle selector with PCD, it gains
$+4.4^*$ mIoU and $+4.0^*$ GQ@0.5. Knowing which tokens to keep does not stop the prior, and keeping only the gold
frames' tokens starves the rest of the timeline. The strongest oracle sees only the evidence and also uses PCD.
Against it, \method is level on grounding ($-2.2$ mIoU and $-1.5$ GQ@0.5, both n.s.) but answers worse ($-5.3^*$
accuracy). Even the gold window leaks 45\%, so with every frame inside the evidence the model still often starts near
0\,s. On the gold window, PCD lowers the leak to 35\% and adds $+5.5^*$ mIoU, but it does not change the answers
($-2.3$, n.s.). For Qwen3-VL-8B, then, more evidence raises accuracy, while PCD fixes the time values.

For TimeLens-8B (\cref{tab:oracle-tl8b}), the bounds are much higher and \method is further from them. Over the
oracle selector with PCD, \method is level on grounding, with $+0.5$ mIoU and $+2.7$ GQ@0.5 (n.s.). The accuracy
difference is $+4.8$, but its interval reaches down to 0.0. The gold window alone is $+13.2^*$ mIoU above all tokens (both at
batch size 1). Adding PCD on the gold window raises mIoU by $+8.3^*$, and it also raises accuracy ($+8.0^*$) and
GQ@0.5 ($+6.7^*$). \method stays below this bound by $11.0^*$ mIoU, $9.7^*$ accuracy and $8.8^*$ GQ@0.5. So for the
grounding-tuned model, PCD and evidence work together. With the right frames, contrasting the time values also
improves the answers. The gap to the gold window shows how much a better choice of frames could still add for this
model.

%% ================================================================================================= 6.8 retrieval
\subsection{Analysis of BOLT}
\label{sec:res-retrieval}

The SigLIP~2 index does find the evidence. Over the seconds before each question, the AUROC of evidence seconds
against the rest is 0.663. The top-ranked second lies in the gold interval for 63\% of questions (chance 36\%), and
for 41\% of location questions (chance 13.5\%). BOLT frames put 40\% of the memory inside the evidence,
against 36\% for uniform frames. This helps random tokens ($+3.8^*$ mIoU). Once tokens are chosen by the question,
however, BOLT adds only $+0.8$ to $+0.9$ (n.s.), because both stages focus the memory on the question. Subtracting a
question-free prior from SigLIP's score does not help (AUROC 0.635--0.666). This is expected for a retriever trained
with a contrastive loss, whose score already estimates pointwise mutual information~\cite{siglip2}. Re-weighting the
token budget of uniform frames by the same relevance adds $+1.3$ (n.s.). The choice of encoded frames matters more
than how the budget is split among fixed frames.

%% ================================================================================================= 6.9 negative
\subsection{What did not work}
\label{sec:negative}

\begin{table}[t]
  \centering
  \small
  \setlength{\tabcolsep}{3pt}
  \caption{\textbf{Negative results} (S-EMBER; $\Delta$mIoU against the stated reference).}
  \label{tab:negative}
  \begin{adjustbox}{max width=\linewidth}\begin{tabular}{@{}lll@{}}
    \toprule
    Idea & Reference & $\Delta$mIoU \\
    \midrule
    VCD-style noised-frame contrast   & greedy          & $+0.5$ \\
    Permuted-timestamp contrast       & greedy          & $+0.9$ \\
    Evidence-against decoding rule    & PCD             & \dn{$-5.6^*$} \\
    Attention gain while decoding     & greedy          & $+1.3$ \\
    Bias stated in the instruction    & greedy          & $-0.1$ \\
    Relevance-guided frame zoom       & greedy          & $-1.5$ \\
    Zoom inside the gold window       & greedy          & $+2.3$ \\
    Unequal per-frame quotas (5 variants) & equal quota & $-0.4$ to $+0.5$ \\
    Contrastive token score           & question top-$k$ & $-1.2$ to $+0.5$ \\
    Prior-subtracted retrieval        & SigLIP          & AUROC $\leq +0.003$ \\
    128 frames at the same budget     & 64 frames       & $+1.2$ / \dn{$-3.6^*$} \\
    \bottomrule
  \end{tabular}\end{adjustbox}
\end{table}

\Cref{tab:negative} collects the ideas that gave little or nothing. Taken together, they suggest that at a fixed
budget, redistributing evidence by zooming, re-weighting frames or unequal quotas is close to its ceiling. Large gains
came only from removing the prior where it acts (PCD), or from selecting tokens by the actual question in every
layer, with or without the equal quota (\cref{sec:res-coverage}).

%% ================================================================================================= 6.10 errors
\subsection{Where answers still fail}
\label{sec:errors}

\begin{table}[t]
  \centering
  \small
  \setlength{\tabcolsep}{3pt}
  \caption{\textbf{Error types} on S-EMBER, pooled over eight memories, greedy \vs PCD
  (earlier implementation of PCD).}
  \label{tab:errors}
  \begin{adjustbox}{max width=\linewidth}\begin{tabular}{@{}lccc@{}}
    \toprule
    Outcome & Greedy & PCD & $\Delta$ \\
    \midrule
    IoU $\geq 0.5$, answer correct      & \phantom{0}6.1\% & \phantom{0}9.9\% & \up{$+3.8^*$} \\
    IoU $\geq 0.5$, answer wrong        & 16.3\% & 20.7\% & \up{$+4.4^*$} \\
    Partial overlap (IoU $<0.5$)        & 17.0\% & 15.7\% & $-1.3$ \\
    Inside gold, too short              & 22.2\% & 18.5\% & \dn{$-3.7^*$} \\
    Entirely \emph{before} gold         & 33.1\% & 18.9\% & $-14.2^*$ \\
    Entirely after gold                 & \phantom{0}2.7\% & \phantom{0}5.7\% & \up{$+3.0^*$} \\
    Covers gold, too long               & \phantom{0}2.2\% & \phantom{0}2.5\% & $+0.3$ \\
    No answer / unparsed           & \phantom{0}0.5\% & \phantom{0}8.1\% & \up{$+7.6^*$} \\
    \bottomrule
  \end{tabular}\end{adjustbox}
\end{table}

\Cref{tab:errors} shows that PCD works where it was aimed. Intervals entirely before the gold one are the signature
of the prior. Under greedy decoding they are the largest error type, and with PCD they fall to less than a fifth of
answers. The remaining failures point to four limits.
(i) Wrong answers on good intervals. Twice as many answers are well grounded but wrong as are grounded and right. Counting answers do not use
the events the model itself listed, and an IoU of 0.5 still allows a two-fold duration error.
(ii) Recent events. Of the intervals that still end before the gold one, 81\% have evidence in the second half of the
window, and the model picks an earlier similar event instead. When the evidence lies in the last third of the window,
GQ@0.5 is 4.4.
(iii) Sparse evidence. With fewer than 4 of the 64 frames inside the evidence, GQ@0.5 is 1--2\% for any method,
against 15\% with 16 or more. No decoding rule can recover evidence that is not in the input.
(iv) Label ambiguity. On the 98 questions whose annotators disagree (mean pairwise IoU $<0.6$), every method
scores near zero.
